\section{Total-variation estimation and interpretation}

The signed experiment in \cref{obj:prop:signed-causal-bridge} connects the welfare analysis to estimation of total-variation distance from a public uniform reference. The resulting statistical characterization concerns a family with equal mass across pairs and an unknown signed contrast within each pair. A separate converse applies to the full probability simplex on the same alphabet. To state these conclusions, we first specify the two statistical models and their common estimand.

\begin{definitionv}{def:tv-corollary-models}[Total-variation models $\mathcal P_d^{\mathrm{pair}}$ and $\Delta_{2d}$]
Define the finite paired alphabet and its probability simplex by
\[
\mathcal A_d=[d]\times\{-1,1\},
\qquad
\Delta_{2d}
=
\left\{
\mathsf p:\mathcal A_d\to[0,1]:
\sum_{j=1}^d\sum_{s\in\{-1,1\}}\mathsf p(j,s)=1
\right\}.
\]
Here \(\mathsf p\) is a probability law on the paired alphabet. Define the paired family and its uniform-reference total-variation functional by
\[
\mathcal P_d^{\mathrm{pair}}
=
\{P_\theta:\theta\in[-1/2,1/2]^d\}
\subseteq\Delta_{2d},
\]
\[
D_{\mathrm{TV}}(\mathsf p)
=
\frac12\sum_{j=1}^d\sum_{s\in\{-1,1\}}
\left|\mathsf p(j,s)-\frac1{2d}\right|.
\]
The reference law \(\mathsf U_{2d}\) is public and uniform on \(\mathcal A_d\). In either statistical model
\[
\mathcal M=\mathcal P_d^{\mathrm{pair}}
\quad\text{or}\quad
\mathcal M=\Delta_{2d},
\]
the data are \(n\) independent, identically distributed draws from one unknown law \(\mathsf p\in\mathcal M\), and the estimand is \(D_{\mathrm{TV}}(\mathsf p)\). The causal model \(\mathcal V_d\) retains its full-data laws and welfare estimand \(V(P)\).
\end{definitionv}

The paired family \leanref{sym:\mathcal P_d^{\mathrm{pair}}}{\ensuremath{\mathcal P_d^{\mathrm{pair}}}} fixes each pair's total probability at \(1/d\), whereas the simplex \leanref{sym:\Delta_{2d}}{\ensuremath{\Delta_{2d}}} permits arbitrary probability masses on the alphabet. In both models, the scalar target \leanref{sym:D_{\mathrm{TV}}}{\ensuremath{D_{\mathrm{TV}}(\mathsf p)}} measures departure from the known reference law \leanref{sym:\mathsf U_{2d}}{\ensuremath{\mathsf U_{2d}}}. This common reference makes the model comparison precise: the input alphabet and target remain fixed while the set of possible sampling laws expands.

The estimation and interval criteria use pure local privacy across every pair of input symbols. The following specification also fixes the information available to the analyst and the meaning of sequential interaction.

\begin{definitionv}{def:tv-private-decision-values}[Private total-variation decision values $\mathfrak R^{\mathrm{TV}}_{\mathcal C}$ and $\mathfrak H^{\mathrm{TV}}_{\mathcal C}$]
The class \(\mathfrak K_{\mathcal C}(n,d,\varepsilon)\), for \(\mathcal C\in\{\mathrm{NI},\mathrm{SI}\}\), consists of one-message-per-participant protocols on the paired input alphabet \(\mathcal A_d=[d]\times\{-1,1\}\). A protocol \(\mathsf K\) satisfies the following conditions:
\begin{itemize}
\item \textbf{(Spaces and kernels.)} The protocol has an arbitrary standard-Borel public-seed space with probability law \(\lambda\), standard-Borel participant message spaces, and participant-specific message kernels \(\mathsf K_i\) conditional on the input, preceding message history, and public seed.
\item \textbf{(Local privacy.)} For every participant \(i\), public seed \(r\), preceding message history \(\eta\), Borel message event \(E\), \(j,j'\in[d]\), and \(s,s'\in\{-1,1\}\),
\[
\mathsf K_i(E\mid(j,s),\eta,r)
\le e^\varepsilon\mathsf K_i(E\mid(j',s'),\eta,r).
\]
\item \textbf{(Public specification.)} All protocol choices depend only on \(n,d,\varepsilon\), the specified model, and public bounds.
\item \textbf{(Interaction.)} For \(\mathcal C=\mathrm{NI}\), every kernel is constant in its message history; participant-specific kernels and shared public seeds are permitted. For \(\mathcal C=\mathrm{SI}\), kernels may depend on preceding messages, and each participant is accessed once.
\item \textbf{(Independent randomizers.)} The public seed \(R\) and analyst randomizer \(U\) have laws
\[
R\sim\lambda,
\qquad
U\sim\operatorname{Uniform}[0,1],
\]
and are mutually independent and independent of the independent, identically distributed inputs.
\end{itemize}

Write
\[
\mathbf Z^{\mathsf K}=(Z_i^{\mathsf K})_{i=1}^n
\]
for the resulting message vector. An estimate \(\widehat D\) is any Borel real-valued function of \((\mathbf Z^{\mathsf K},R,U)\). An interval \(J_{\mathrm{TV}}\) is any nonempty connected closed interval contained in \([0,1]\), with ordered Borel endpoints that depend only on the same information.

For either statistical model
\[
\mathcal M\in\{\mathcal P_d^{\mathrm{pair}},\Delta_{2d}\},
\]
the minimax squared-error risk and globally honest expected interval length for the uniform-reference total-variation distance \(D_{\mathrm{TV}}(\mathsf p)\) are
\[
\begin{aligned}
\mathfrak R^{\mathrm{TV}}_{\mathcal C}(\mathcal M;n,d,\varepsilon)
&=
\inf_{\mathsf K\in\mathfrak K_{\mathcal C}(n,d,\varepsilon)}
\inf_{\widehat D}
\sup_{\mathsf p\in\mathcal M}
\mathbb E_{\mathsf p,\mathsf K}
\bigl[(\widehat D-D_{\mathrm{TV}}(\mathsf p))^2\bigr],\\
\mathfrak H^{\mathrm{TV}}_{\mathcal C}(\mathcal M;n,d,\varepsilon)
&=
\inf_{\mathsf K\in\mathfrak K_{\mathcal C}(n,d,\varepsilon)}
\inf_{\substack{J_{\mathrm{TV}}:\,
\inf_{\mathsf p\in\mathcal M}
\Pr_{\mathsf p,\mathsf K}
\{D_{\mathrm{TV}}(\mathsf p)\in J_{\mathrm{TV}}\}\ge0.90}}
\sup_{\mathsf p\in\mathcal M}
\mathbb E_{\mathsf p,\mathsf K}
\operatorname{length}J_{\mathrm{TV}}.
\end{aligned}
\]
Here \(\operatorname{length}J_{\mathrm{TV}}\) is the difference between its ordered endpoints. Probabilities and expectations include input sampling, channel randomness, public randomness, and analyst randomization. Singleton intervals are permitted. For the paired model, the attaining intervals specified below are contained in \([0,1/4]\).
\end{definitionv}

Global honesty requires coverage at every sampling law in the selected model, and expected length is evaluated at the worst such law. The connected interval \leanref{sym:J_{\mathrm{TV}}}{\ensuremath{J_{\mathrm{TV}}}} therefore measures the precision of a single reported range for the scalar target. The noninteractive and sequential classes both allow shared public randomness and participant-specific message distributions; their distinction is whether a participant's message distribution may depend on preceding messages.

\subsection{Matching precision on the paired family}

For the paired family, total variation is one half of the average absolute contrast in the signed experiment. This identity carries the approximation and private-moment analysis into the statistical problem. The published Bernstein approximation limit used in the comparison is given by \citet[Section 3.1, p.~8]{CaiLow2011Nonsmooth}. Under this input, the following result supplies matching squared-error and honest-length bounds, together with a concrete noninteractive procedure attaining their orders.

\begin{theoremv}{thm:paired-uniform-tv-frontier}[Paired-uniform total-variation frontier]*
This result is conditional on the cited, unformalized Bernstein absolute-approximation limit. If \(e_K=\inf_{\deg p\le K}\sup_{|x|\le1}\lvert |x|-p(x)\rvert\) is the best uniform error for approximating \(|x|\) on \([-1,1]\), then \(2k e_{2k}\to\beta_*\) for a constant \(0.278<\beta_*<0.286\). See \citet[Section 3.1]{CaiLow2011Nonsmooth} and \citet{Bernstein1913}.

For the paired family in \cref{obj:def:tv-corollary-models}, every integer \(d\ge2\) and every \(\theta\in[-1/2,1/2]^d\) satisfy
\[
 P_\theta(j,s)=\frac{1+s\theta_j}{2d},
 \qquad
 D_{\mathrm{TV}}(P_\theta)
 =\operatorname{TV}(P_\theta,\mathsf U_{2d})
 =\frac12F(\theta)
 =\frac{1}{2d}\sum_{j=1}^{d}|\theta_j|.
\]
Here \(\mathsf U_{2d}\) is the uniform law on the paired alphabet and \(F(\theta)\) is the average absolute contrast.

There are universal real constants \(0<c\le C_0<\infty\) such that the following conclusions hold for every resource tuple satisfying
\begin{itemize}
\item \textbf{(Sample size.)} \(n\) is an integer with \(n\ge2\).
\item \textbf{(Dimension.)} \(d\) is an integer with \(d\ge2\).
\item \textbf{(Privacy.)} \(0<\varepsilon\le1\).
\end{itemize}
For either \(\mathcal C\in\{\mathrm{NI},\mathrm{SI}\}\), let
\[
 t=n\varepsilon^2,\qquad
 r_{\mathcal C}(n,d,\varepsilon)=\rho(t,d),\qquad
 h_{\mathcal C}(n,d,\varepsilon)=\sqrt{\rho(t,d)},
\]
where \(\rho\) is the evaluated rate defined in \cref{obj:def:frontier-handle}. The decision values of \cref{obj:def:tv-private-decision-values} obey
\[
\begin{aligned}
 c\,r_{\mathcal C}(n,d,\varepsilon)
 &\le
 \mathfrak R^{\mathrm{TV}}_{\mathcal C}
   (\mathcal P_d^{\mathrm{pair}};n,d,\varepsilon)
 \le C_0\,r_{\mathcal C}(n,d,\varepsilon),\\
 c\,h_{\mathcal C}(n,d,\varepsilon)
 &\le
 \mathfrak H^{\mathrm{TV}}_{\mathcal C}
   (\mathcal P_d^{\mathrm{pair}};n,d,\varepsilon)
 \le C_0\,h_{\mathcal C}(n,d,\varepsilon).
\end{aligned}
\]

For each such resource tuple there exists a pure \(\varepsilon\)-locally private noninteractive paired-input protocol \(\mathsf K\), with each participant's message space in bijection with \(\{-1,1\}^d\), together with a transcript-only estimator \(\widehat D\) and a connected closed interval \(J_{\mathrm{TV}}\subseteq[0,1/4]\), such that, for every \(\mathsf p\in\mathcal P_d^{\mathrm{pair}}\), under independent sampling from \(\mathsf p\),
\[
\begin{aligned}
 \mathbb E_{\mathsf p,\mathsf K}
   [(\widehat D-D_{\mathrm{TV}}(\mathsf p))^2]
 &\le C_0\,r_{\mathcal C}(n,d,\varepsilon),\\
 \mathbb P_{\mathsf p,\mathsf K}
   \{D_{\mathrm{TV}}(\mathsf p)\in J_{\mathrm{TV}}\}
 &\ge0.90,\\
 \mathbb E_{\mathsf p,\mathsf K}
   [\operatorname{length}(J_{\mathrm{TV}})]
 &\le C_0\,h_{\mathcal C}(n,d,\varepsilon).
\end{aligned}
\]
These guarantees also hold for the concrete resource-selected sign-vector polynomial protocol and its total-variation estimator and interval, with the construction specified in \cref{obj:def:frontier-handle,obj:thm:uniform-private-value-frontiers}: this protocol is noninteractive and pure \(\varepsilon\)-locally private, and its squared risk, coverage, and expected length satisfy the three displayed bounds with \(\mathcal C=\mathrm{NI}\).

For the same universal constants and every admissible resource tuple,
\[
\begin{aligned}
1
&\le
\frac{\mathfrak R^{\mathrm{TV}}_{\mathrm{NI}}
  (\mathcal P_d^{\mathrm{pair}};n,d,\varepsilon)}
 {\mathfrak R^{\mathrm{TV}}_{\mathrm{SI}}
  (\mathcal P_d^{\mathrm{pair}};n,d,\varepsilon)}
\le \frac{C_0}{c},\\
1
&\le
\frac{\mathfrak H^{\mathrm{TV}}_{\mathrm{NI}}
  (\mathcal P_d^{\mathrm{pair}};n,d,\varepsilon)}
 {\mathfrak H^{\mathrm{TV}}_{\mathrm{SI}}
  (\mathcal P_d^{\mathrm{pair}};n,d,\varepsilon)}
\le \frac{C_0}{c}.
\end{aligned}
\]

The rate has the following exact resource properties. Write \(L=\log(ed)\). For every integer \(d\ge2\) and every \(0<t<d^2L\),
\[
 \rho(t,d)=1
 \quad\Longleftrightarrow\quad
 t\le\frac{d^2L}{ed-e}.
\]
For every sequence of admissible resource tuples
\((n_k,d_k,\varepsilon_k)\), write
\[
 t_k=n_k\varepsilon_k^2,\qquad L_k=\log(ed_k).
\]
Then
\[
 \rho(t_k,d_k)\longrightarrow0
 \quad\Longleftrightarrow\quad
 t_k\longrightarrow\infty
 \ \text{and}\
 \frac{\log(1+d_k^2/t_k)}{L_k}\longrightarrow0.
\]
Moreover,
\[
 \rho(t_k,d_k)\asymp t_k^{-1}
 \quad\Longleftrightarrow\quad
 \bigl(\exists a>0:\ t_k\ge a\ \text{eventually}\bigr)
 \ \text{and}\
 \bigl(\exists M\in\mathbb R:\ \min\{t_k,d_k\}\le M\
       \text{eventually}\bigr).
\]
Here \(\asymp\) means two-sided comparison by fixed positive finite constants for all sufficiently large \(k\). If \(t_k\to\infty\), this equivalence becomes
\[
 \rho(t_k,d_k)\asymp t_k^{-1}
 \quad\Longleftrightarrow\quad
 \exists M\in\mathbb N:\ d_k\le M\ \text{eventually}.
\]
For every pair of real numbers \(0<a\le b\), there exist real constants \(0<c\le C_0<\infty\) and an integer \(D_0\) such that, for every integer \(d\ge D_0\) with \(d\ge2\) and every real \(t\) satisfying \(a\le t/d^2\le b\),
\[
 c\left(\frac{\log L}{L}\right)^2
 \le \rho(t,d)
 \le C_0\left(\frac{\log L}{L}\right)^2.
\]

Finally, these sequence properties transfer to the actual minimax values for each fixed \(\mathcal C\in\{\mathrm{NI},\mathrm{SI}\}\). Along every admissible sequence,
\[
\begin{aligned}
 \mathfrak R^{\mathrm{TV}}_{\mathcal C}
   (\mathcal P_{d_k}^{\mathrm{pair}};n_k,d_k,\varepsilon_k)
   \longrightarrow0
 &\quad\Longleftrightarrow\quad
 t_k\longrightarrow\infty
 \ \text{and}\\
 \frac{\log(1+d_k^2/t_k)}{L_k}\longrightarrow0,\\
 \mathfrak H^{\mathrm{TV}}_{\mathcal C}
   (\mathcal P_{d_k}^{\mathrm{pair}};n_k,d_k,\varepsilon_k)
   \longrightarrow0
 &\quad\Longleftrightarrow\quad
 t_k\longrightarrow\infty
 \ \text{and}\\
 \frac{\log(1+d_k^2/t_k)}{L_k}\longrightarrow0.
\end{aligned}
\]
The squared-error value satisfies
\[
 \mathfrak R^{\mathrm{TV}}_{\mathcal C}
   (\mathcal P_{d_k}^{\mathrm{pair}};n_k,d_k,\varepsilon_k)
   \asymp t_k^{-1}
 \quad\Longleftrightarrow\quad
 \bigl(\exists a>0:\ t_k\ge a\ \text{eventually}\bigr)
 \ \text{and}\
 \bigl(\exists M\in\mathbb R:\ \min\{t_k,d_k\}\le M\
       \text{eventually}\bigr).
\]
Under \(t_k\to\infty\), equivalently,
\[
 \mathfrak R^{\mathrm{TV}}_{\mathcal C}
   (\mathcal P_{d_k}^{\mathrm{pair}};n_k,d_k,\varepsilon_k)
   \asymp t_k^{-1}
 \quad\Longleftrightarrow\quad
 \exists M\in\mathbb N:\ d_k\le M\ \text{eventually}.
\]
For each fixed \(\mathcal C\) and every \(0<a\le b\), there exist real constants \(0<c\le C_0<\infty\) and an integer \(D_0\) such that every admissible tuple with \(d\ge D_0\) and \(a\le n\varepsilon^2/d^2\le b\) satisfies
\[
\begin{aligned}
 c\left(\frac{\log L}{L}\right)^2
 &\le
 \mathfrak R^{\mathrm{TV}}_{\mathcal C}
   (\mathcal P_d^{\mathrm{pair}};n,d,\varepsilon)
 \le C_0\left(\frac{\log L}{L}\right)^2,\\
 \sqrt c\left|\frac{\log L}{L}\right|
 &\le
 \mathfrak H^{\mathrm{TV}}_{\mathcal C}
   (\mathcal P_d^{\mathrm{pair}};n,d,\varepsilon)
 \le \sqrt{C_0}\left|\frac{\log L}{L}\right|.
\end{aligned}
\]
\end{theoremv}
% CAUSALSMITH-CITED-SCOPE-BEGIN thm:paired-uniform-tv-frontier
\verificationfootnotetext{\textbf{Formalization scope.} This result uses the cited conclusion \citep{CaiLow2011Nonsmooth} (Section 3.1, definition of the Bernstein constant and immediately following paragraph, page 8). The cited source proof is not formalized here: Lean verifies its use and all remaining steps. If that cited conclusion is false, the portions of this result that depend on it are not certified.}
% CAUSALSMITH-CITED-SCOPE-END thm:paired-uniform-tv-frontier

The identity in \cref{obj:thm:paired-uniform-tv-frontier} explains the common precision benchmark: departure from uniformity within a pair contributes the absolute value of its signed contrast. Thus the scalar total-variation target retains the same absolute-value structure that governs welfare in \cref{obj:prop:signed-causal-bridge}. The concrete paired protocol, estimator, and interval in \cref{obj:synth_5} use the same private-moment and polynomial-approximation ingredients as the welfare construction in \cref{obj:def:frontier-handle}. They attain both the squared-error order and the square-root order for globally honest expected interval length.

The resource characterizations describe when these guarantees yield increasing precision. Both decision values vanish exactly when effective private resource diverges and the normalized logarithmic resource deficit vanishes. Along sequences with diverging effective private resource, the usual inverse-resource squared-error order is equivalent to an eventually bounded dimension. When effective private resource is proportional to the square of the dimension, the theorem instead identifies squared-error order \((\log L/L)^2\) and honest-length order \(|\log L/L|\), where \(L=\log(ed)\) is the log-dimension factor.

The interaction comparison is uniform over admissible sample sizes, dimensions, and privacy budgets. Within the one-person-one-message classes of \cref{obj:def:tv-private-decision-values}, sequential adaptation can improve the minimax criteria by at most a universal multiplicative factor. A noninteractive procedure already attains their common orders. This comparison concerns adaptation to preceding participants' messages under the same pure local privacy requirement and with each participant accessed once.

\subsection{A converse on the full simplex}

The paired family is a subset of the full simplex in \cref{obj:def:tv-corollary-models}. Consequently, estimation over the simplex includes the paired-family problem, and global simplex coverage includes coverage over the paired family. With the same alphabet and protocol restrictions, this inclusion transfers the paired lower bounds to arbitrary probability laws. The approximation input remains the Bernstein limit cited in \citet[Section 3.1, p.~8]{CaiLow2011Nonsmooth}.

\begin{theoremv}{thm:full-simplex-tv-converse}[Full-simplex total-variation lower bounds]*
This result is conditional on the cited, unformalized Bernstein absolute-approximation limit. If \(e_K=\inf_{\deg p\le K}\sup_{|x|\le1}\lvert |x|-p(x)\rvert\) is the best uniform error for approximating \(|x|\) on \([-1,1]\), then \(2k e_{2k}\to\beta_*\) for a constant \(0.278<\beta_*<0.286\). See \citet[Section 3.1]{CaiLow2011Nonsmooth} and \citet{Bernstein1913}.

Let \(\Delta_{2d}\) be the probability simplex on the paired alphabet \(\mathcal A_d=[d]\times\{-1,1\}\), and let
\[
D_{\mathrm{TV}}(\mathsf p)=\operatorname{TV}(\mathsf p,\mathsf U_{2d}),
\]
where \(\mathsf U_{2d}\) is the uniform reference law. There is a universal constant \(c>0\) such that, for every sample size \(n\), dimension \(d\), and privacy budget \(\varepsilon\) satisfying
\begin{itemize}
\item \textbf{(Sample size.)} \(n\ge2\).
\item \textbf{(Dimension.)} \(d\ge2\).
\item \textbf{(Privacy budget.)} \(0<\varepsilon\le1\).
\end{itemize}
and every protocol-class label \(\mathcal C\in\{\mathrm{NI},\mathrm{SI}\}\), the private decision values of \cref{obj:def:tv-private-decision-values}, with iid inputs drawn from \(\mathsf p\in\Delta_{2d}\), satisfy
\[
\begin{aligned}
\mathfrak R^{\mathrm{TV}}_{\mathcal C}(\Delta_{2d};n,d,\varepsilon)
&\ge c\,r_{\mathcal C}(n,d,\varepsilon),\\
\mathfrak H^{\mathrm{TV}}_{\mathcal C}(\Delta_{2d};n,d,\varepsilon)
&\ge c\,h_{\mathcal C}(n,d,\varepsilon),
\end{aligned}
\]
where the squared-error and connected-length rates are
\[
r_{\mathcal C}(n,d,\varepsilon)=\rho(n\varepsilon^2,d),
\qquad
h_{\mathcal C}(n,d,\varepsilon)=\sqrt{r_{\mathcal C}(n,d,\varepsilon)}.
\]
The interval value uses connected closed intervals in \([0,1]\) with coverage at least \(0.90\) at every law in \(\Delta_{2d}\). The protocol classes permit arbitrary public seeds, participant-specific kernels, standard-Borel message spaces, and, for the sequential class, dependence on preceding messages, with one message per participant and pure local privacy.
\end{theoremv}
% CAUSALSMITH-CITED-SCOPE-BEGIN thm:full-simplex-tv-converse
\verificationfootnotetext{\textbf{Formalization scope.} This result uses the cited conclusion \citep{CaiLow2011Nonsmooth} (Section 3.1, definition of the Bernstein constant and immediately following paragraph, page 8). The cited source proof is not formalized here: Lean verifies its use and all remaining steps. If that cited conclusion is false, the portions of this result that depend on it are not certified.}
% CAUSALSMITH-CITED-SCOPE-END thm:full-simplex-tv-converse

The intuition behind \cref{obj:thm:full-simplex-tv-converse} is that a larger model retains the difficult sampling laws contained in its paired subfamily. Allowing arbitrary pair masses therefore preserves the lower bounds for the minimax squared-error value \leanref{sym:\mathfrak R^{\mathrm{TV}}}{\ensuremath{\mathfrak R^{\mathrm{TV}}_{\mathcal C}}} and the globally honest expected-length value \leanref{sym:\mathfrak H^{\mathrm{TV}}}{\ensuremath{\mathfrak H^{\mathrm{TV}}_{\mathcal C}}}. The result supplies necessary precision bounds on the simplex, while \cref{obj:thm:paired-uniform-tv-frontier} supplies a matching characterization on the paired family.

\subsection{Interpretation for welfare reporting}

The causal result in \cref{obj:thm:uniform-private-value-frontiers} gives a precision benchmark for reporting a fixed population's first-best welfare under the model in \cref{obj:def:causal-model}. The population law determines the target, and independent participant sampling and private messages determine the uncertainty in its estimate. Squared-error risk measures the accuracy of the reported welfare value; global honesty requires coverage across the declared population laws, with expected connected interval length measuring reporting precision.

The signed bridge in \cref{obj:prop:signed-causal-bridge} explains why the paired statistical problem is informative for this interpretation. Both analyses involve the average absolute stratum contrast, and \cref{obj:thm:paired-uniform-tv-frontier} isolates that component as a distance from a known uniform law. The welfare and paired-family interaction comparisons in \cref{obj:thm:uniform-private-value-frontiers,obj:thm:paired-uniform-tv-frontier} establish common minimax orders for noninteractive and sequential protocols within their respective declared classes. These conclusions give a sample-size and privacy benchmark for scalar reporting under independent sampling and one message per participant.

\subsection{Limitations and future work}

The causal model uses public uniform stratum probabilities and fair treatment assignment. Unknown stratum probabilities or treatment propensities would require an analysis of their estimation and their contribution to welfare precision under local privacy. Repeated records per participant would likewise require a participant-level privacy specification and a sampling model that accounts for dependence within participants. The present interaction comparison concerns protocols that access each participant once.

The displayed procedure is a conservative witness to uniform rate attainment rather than a calibrated recommendation for applied reporting. Its polynomial branches require \(L=\log(ed)\ge4096\), hence \(d\ge e^{4095}\), and each active participant releases a \(d\)-coordinate sign vector. Its welfare interval uses radius \(\sqrt{10^{17}\rho(t,d)}\); whenever \(\rho(t,d)\ge2.5\times10^{-18}\), clipping makes that interval span the entire welfare range regardless of its center. Sharper constants and communication-efficient procedures are therefore important practical questions even though these conservative choices establish the stated uniform rates.

The matching bounds characterize minimax orders up to universal constants. Sharp leading constants and procedures optimized for those constants remain topics for further work. Learned-policy regret is a separate target from estimation of the population's first-best welfare: it would require evaluating the welfare achieved by a policy selected from private data. Design-based inference for a finite population would require uncertainty and coverage statements based on the treatment-assignment mechanism with the population held fixed, rather than the independent population sampling used here.

Finally, \cref{obj:thm:full-simplex-tv-converse} supplies lower bounds for total-variation estimation and honest intervals on the full simplex. Matching full-simplex upper bounds remain an open question; the attaining guarantees in \cref{obj:thm:paired-uniform-tv-frontier} apply to the paired family.
